\chapter{Berechnung von $U_C$ bei schwacher Dotierung}\label{sec:anhang_b}

Es soll gezeigt werden, dass für alle $\mu \neq 0$, für die die Näherung $\rho(\epsilon) \approx A \abs{\epsilon}$ gilt, die kritische Kopplung $U_C$ verschwindet.

Dazu betrachte die Selbstkonsistenzgleichung für beliebige Füllung~\eqref{eq:gap_eq_arbitrary_filling_integral} bei $T= 0$ und für $\Delta_0 \neq 0$:
\begin{align}
1 = \frac{UA}{2} \int_{\mu - E_D}^{\mu + E_D} \frac{\abs{\epsilon}}{\sqrt{(\epsilon - \mu)^2 + \Delta_0^2}} \, d\epsilon = \frac{UA}{2} \int_{-E_D}^{E_D} \frac{\abs{\xi + \mu}}{\sqrt{\xi^2 + \Delta_0^2}} \,  d\xi
\end{align}
mit $\xi = \epsilon - \mu$. Es gilt also
\begin{align}
    U(\Delta_0) = \frac{2}{A\cdot F(\Delta_0)}
\end{align}
mit
\begin{align}
    F(\Delta_0) &= \int_{-E_D}^{E_D} \frac{\abs{\xi + \mu}}{\sqrt{\xi^2 + \Delta_0^2}} \,  d\xi \, .
\end{align}
Da die kritische Kopplung $U_C$ über den Limes $\Delta_0 \rightarrow 0$ von $U$ definiert ist, verschwindet $U_C$ genau dann, wenn $F(\Delta_0)$ divergiert. Zur Auswertung des Betrags im Integranden werden zwei Fälle unterschieden.

\subsubsection{Fall $\abs{\mu} < E_D$}

Das Integral wird per Fallunterscheidung aufgeteilt:
\begin{align}
    F(\Delta_0) &= -\int_{-E_D}^{-\mu} \frac{\xi + \mu}{\sqrt{\xi^2 + \Delta_0^2}} \, d\xi + \int_{-\mu}^{E_D} \frac{\xi + \mu}{\sqrt{\xi^2 + \Delta_0^2}}\, d\xi = g(-E_D) + g(E_D) \, ,
\end{align}
und mit
\begin{align}
    g(x) = \int_{-\mu}^{x} \frac{\xi + \mu}{\sqrt{\xi^2 + \Delta_0^2}}\, d\xi = \sqrt{x^2 + \Delta_0^2} - \sqrt{\mu^2 + \Delta_0^2} + \mu \left(\operatorname{arsinh}\left(\frac{x}{\Delta_0}\right) - \operatorname{arsinh}\left(\frac{-\mu}{\Delta_0}\right)\right)
\end{align}
ergibt sich
\begin{align}
    F(\Delta_0) &= 2\left(\sqrt{E_D^2 + \Delta_0^2} - \sqrt{\mu^2 + \Delta_0^2}\right) + \mu \left[\operatorname{arsinh}\left(\frac{E_D}{\Delta_0}\right) + \operatorname{arsinh}\left(\frac{-E_D}{\Delta_0}\right) - 2 \operatorname{arsinh}\left(\frac{-\mu}{\Delta_0}\right)  \right] \, .
\end{align}
Die Wurzelterme in $F(\Delta_0)$ konvergieren einfach gegen eine Konstante, weswegen also nur noch die Konvergenz des $\arsinh$-Terms überprüft werden muss.

Aufgrund der Punktsymmetrie von $\arsinh$ kann dieser jedoch stark vereinfacht werden, da somit $\arsinh(E_D/\Delta_0) + \arsinh(-E_D/\Delta_0) = 0$ gilt und nur $-2\mu \arsinh(-\mu/\Delta_0)$ übrig bleibt. Für $\Delta_0 \rightarrow 0$ divergiert das Argument und der übrigbleibende $\arsinh$ divergiert daher ebenfalls logarithmisch.

Für $\mu = 0$ entfällt diese logarithmische Divergenz und $F(\Delta_0) \rightarrow F(0) = 2E_D$, man erhält also das Ergebnis aus~\eqref{eq:U_C}.

\subsubsection{Fall $\abs{\mu} \geq E_D$}

In diesem Fall ändert $\xi + \mu$ im Integrationsintervall $[-E_D, E_D]$ das Vorzeichen nicht mehr, sodass nur noch das Integral
\begin{align}
    F(\Delta_0) = \operatorname{sgn}(\mu) \int_{-E_D}^{E_D} \frac{\xi + \mu}{\sqrt{\xi^2 + \Delta_0^2}}\, d\xi
\end{align}
gelöst werden muss. Analog zum vorherigen Fall ergibt sich
\begin{align}
    F(\Delta_0) = \abs{\mu} \left(\operatorname{arsinh}\left(\frac{E_D}{\Delta_0}\right) - \operatorname{arsinh}\left(\frac{-E_D}{\Delta_0}\right)\right) = 2\abs{\mu} \left(\arsinh\left(\frac{E_D}{\Delta_0}\right)\right)
\end{align}

was für $\Delta_0 \rightarrow 0$ ebenfalls logarithmisch divergiert, solange $\mu \neq 0$ ist.

Damit ist gezeigt, dass die kritische Kopplung, außer bei $\mu = 0$, verschwindet.